\documentclass{article}
\usepackage[T1]{fontenc}
\usepackage{lmodern}
\usepackage{graphicx}
\usepackage{amsmath,amssymb}
\usepackage[a4paper,margin=2cm]{geometry}
\usepackage[absolute,overlay]{textpos}
\setlength{\TPHorizModule}{1mm}
\setlength{\TPVertModule}{1mm}
\usepackage[indonesian]{babel}
\usepackage{hyperref}
\setlength{\emergencystretch}{2em}
% unit: Halaman 6

\begin{document}

% === Halaman 6 (llm) ===

\section*{Perbedaan Kriptografi dan Steganografi}

\begin{itemize}
    \item \textbf{Kriptografi}: menyembunyikan \textit{makna} atau isi (\textit{content}) pesan
    $\rightarrow$ Tujuan: agar pesan tidak dapat dibaca oleh pihak ketiga (lawan)
    
    \item \textbf{Steganografi}: menyembunyikan \textit{keberadaan} (\textit{existence}) pesan
    $\rightarrow$ Tujuan: untuk menghindari kecurigaan (\textit{conspicuous}) dari pihak ketiga (lawan)
\end{itemize}

\end{document}
